\documentclass[twocolumn]{aastex631}
\begin{document}
\title{The reionization ionizing-photon-budget shortfall for star-forming galaxies is not robust to systematics at z\~{}6 (o32 systematic corner)}
\author{NebulaMind Lab (autonomous pipeline)}
\affiliation{NebulaMind Lab, \url{https://lab.nebulamind.net}}
\begin{abstract}
Reionization ionizing-photon-budget reconciliation at z\~{}6: star-forming galaxies require f\_esc=0.048 (+0.048/-0.025) to close the budget (Madau-Dickinson SFRD, log xi\_ion=25.5+/-0.15, clumping C=2-5, JWST-SFRD tail), versus indirect-proxy-inferred f\_esc=0.080 (+0.146/-0.051) from LzLCS O32/beta calibrations. Median delta(required-inferred)=-0.029 dex-frac (16-84\%: -0.170 to +0.034); 34\% of the systematic MC shows a shortfall. Conclusion: the budget CLOSES within the systematic, and the sign holds under both O32 and beta calibrations. The result is bounded by the xi\_ion x clumping x proxy-calibration systematic, not by statistics. Generated autonomously from public data (jwst) via the NebulaMind Lab runner: a bounded, reproducible, descriptive study.
\end{abstract}
\section{Introduction}
Recent studies have highlighted a potential crisis in reconciling the ionizing photon budget during the reionization epoch [Muñoz2024]. This discrepancy arises from the apparent insufficiency of star-forming galaxies to provide enough ionizing photons to account for the observed reionization, given current estimates of their escape fractions and other properties. To address this issue, it is essential to reassess the photon budget using a systematic approach that incorporates published literature values.
\section{Data and method}
In our analysis, we adopt the cosmic star formation rate density (SFRD) from Madau \& Dickinson (2014), along with published calibrations for the ionizing efficiency (xi\_ion) and escape fraction (f\_esc) proxies. Specifically, we use the O32/beta f\_esc proxy calibrations from LzLCS [Chisholm+22, Flury+22] and Simmonds+24. By combining these literature-anchored values, we calculate the ionizing photon budget at z\~{}6 using a systematic method that accounts for uncertainties in xi\_ion, clumping factor (C), and f\_esc.
\section{Result}
Our reconciliation of the reionization ionizing-photon-budget at z\~{}6 reveals that star-forming galaxies require an escape fraction of f\_esc=0.048 (+0.048/-0.025) to close the budget. This value is compared to the indirect-proxy-inferred f\_esc=0.080 (+0.146/-0.051) derived from LzLCS O32/beta calibrations. The median difference between the required and inferred escape fractions is -0.029 dex-frac, with 34\% of systematic Monte Carlo simulations showing a shortfall in the photon budget.
\begin{figure}\centering\includegraphics[width=\columnwidth]{result.png}\caption{The reionization ionizing-photon-budget shortfall for star-forming galaxies is not robust to systematics at z\~{}6 (o32 systematic corner)}\end{figure}
\section{Caveats}
It is crucial to acknowledge that our analysis relies on automated, single-selection, uncalibrated measurements, which may introduce limitations. These include potential biases in the adopted literature values, uncertainties in the calibrations used for xi\_ion and f\_esc proxies, and the assumption of a fixed clumping factor (C) within a narrow range. Furthermore, our approach does not account for other sources of ionizing photons or additional systematic errors that may affect the photon budget calculation. A more comprehensive understanding of these factors is necessary to refine our results and improve the accuracy of reionization models.
\section*{References}
\footnotesize
[Muoz2024] Muñoz et al. (2024). Reionization after JWST: a photon budget crisis?. bibcode 2024MNRAS.535L..37M \par [Davies2021] Davies et al. (2021). The Predicament of Absorption-dominated Reionization: Increased Demands on Ionizing Sources. bibcode 2021ApJ...918L..35D \par [Park2022] Park et al. (2022). Calibrating excursion set reionization models to approximately conserve ionizing photons. bibcode 2022MNRAS.517..192P \par [Duncan2015] Duncan et al. (2015). Powering reionization: assessing the galaxy ionizing photon budget at z < 10. bibcode 2015MNRAS.451.2030D \par [Madau2017] Madau (2017). Cosmic Reionization after Planck and before JWST: An Analytic Approach. bibcode 2017ApJ...851...50M

\end{document}
